\chapter*{Additional Remarks and Development ideas}
\addcontentsline{toc}{chapter}{Evaluation}
\stepcounter{chapter}

This section will outline any possible advancements and considerations that I would like to consider further following submission.

\section{Captive Portal}
A desirable feature for this project is the implementation of a Captive Portal. This is a standard industry practice for public
or guest Wi-Fi networks—such as those found in cafes or cinemas—where users must interact with a landing page before gaining
network access. While this implementation will forgo a full authentication system, a splash page serves as a critical point of
transparency. It allows the user to review terms and conditions and receive notification of potential security risks before proceeding.

The underlying mechanism utilizes DNS hijacking (or DNS redirection), which intercepts outgoing requests and forwards them to
a locally hosted destination. For a more robust deployment, utilities like NoDogSplash can be integrated to handle session management.
This leverages Captive Portal Detection (CPD)—a feature in modern mobile operating systems that automatically triggers a "Sign-In"
prompt when a restricted network is detected, ensuring a seamless user interface.

However, with time constraints this is likely not going to come to fruition.




\section{Camera Placement}
Due to the diverse range of hardware configurations, static camera placement can be challenging; therefore, a dedicated setup
interface is ideal. Providing a live preview allows the user to align themselves correctly, ensuring optimal input for the model.
While transmitting a high-definition video feed introduces significant computational overhead and latency, this can be mitigated
by utilizing a low-resolution or compressed stream during the initialization phase. This replaces guesswork with visual feedback,
significantly improving the user onboarding experience.




\section{Recording}
If users can review video segments and provide manual corrections, these labels can be used to fine-tune the model to a 
specific person. This \textbf{human-in-the-loop} approach could help to tailour the model and improve accuracy.




\section{Differing Models}
Other model types may allow for better accuracy over time. \textbf{Recurrent models} feed data back in with each pass, this could
result in certain aspects influencing future predictions. ne example of this is temporal mapping: as
the model processes a sequence, general trends towards a classification could be aggregated to result in a more consistent pattern
of classfication rather than flickering between states.

Further dimensionality could also be used, considering video input rather than static frames may also be another direction, where
temporal context and a series of frames may be used to inform a more definitive assesment of someones \textbf{tiredness}. For example,
the duration of eye colosure could distinguish between a simple blink and micro-sleep.

Additional resources (such as sound) could also be integrated. While in a static image a user may appear to be yawning, it may be the
case that they are singing, which could otherwise taint the overall classification.



\section{Custom and More Refined Hardware}
The Raspberry Pi 4B, while excellent for prototyping, contains a vast range of unnecessary features for this specific use
case—such as dual-monitor support and high-performance desktop processing—which create unwanted power overhead and increased
unit costs. For a production-ready product, a more tailored hardware stack would be essential to achieve a lower Bill of Materials (BOM).

By moving toward an \textbf{Edge AI} architecture, more discrete and usable hardware could be formed using components such as:
\begin{itemize}
    \item \textbf{ESP32-S3-Sense:} An ultra-low-cost SoC (System on Chip) that includes native AI vector instructions. This allows
    for the running of the model through improved AI optimization, as well as hosting the webpage and managing requests.
    \item \textbf{Kendryte K210:} A dedicated RISC-V AI processor with a built-in KPU (Knowledge Processing Unit). This allows for
    high-speed landmark detection while also reducing power consumption.
    \item \textbf{Custom PCB and Enclosure:} By using a custom casing, the system can be integrated more easily within the vehicle,
    forming a self-contained safety unit. Furthermore, additional lights or buzzers could be added to eliminate the need for external
    peripherals.
    \item \textbf{Infra-Red Camera:} To overcome the limitations in low-light environments that the webcam approach faced, a more
    advanced form of data capture would be appropriate to ensure that the model remains effective during night-time operation.
\end{itemize}
This transition from a general-purpose Single Board Computer (SBC) to a dedicated device that is more compact, boots faster, and
is significantly more affordable for the end consumer would result in a more refined product.